% Slide build of Figure 24.1 -- placement, wrong first.
%
% In print both placements are shown together and the caption says one is an
% outage. On a slide the honest order is the order a student would actually
% arrive at it: put the ACL near the source, because that is the obvious place
% and it saves carrying doomed traffic; watch it take out everything; then move
% it. Being shown the mistake being made is what makes the rule stick, and it
% is the same reasoning behind the book's Break & Fix sections.
%
% The requirement stays on screen at every step, because the whole judgement
% is "does this still meet the requirement" and the room should be able to
% check that without being reminded of it.
\begin{tikzpicture}[font=\scriptsize]
  \useasboundingbox (-1.6,-4.6) rectangle (12.2,4.3);

  \tikzset{
    dv/.style={rounded corners=3pt, draw=#1, fill=#1!8, text=#1!45!black,
               line width=0.9pt, minimum width=1.35cm, minimum height=0.95cm,
               align=center, font=\tiny\bfseries},
    w/.style={line width=0.85pt, neutral!65},
    hdr/.style={font=\scriptsize\bfseries, text=neutral!45!black},
    stepcap/.style={font=\bfseries, anchor=north west, text width=13.4cm,
                    align=left}}

  \node[hdr, text width=13.5cm, align=center] at (5.3,3.5)
      {VLAN 30 must not reach the finance server --- but must still reach
       everything else};

  \node[dv=neutral] (h)   at (0,0.7)    {\faIcon{desktop}\\VLAN 30};
  \node[dv=dom3]    (r1)  at (3.0,0.7)  {\faIcon{route}\\R1};
  \node[dv=dom3]    (r2)  at (7.0,0.7)  {\faIcon{route}\\R2};
  \node[dv=dom1]    (fin) at (10.4,1.75){\faIcon{server}\\Finance};
  \node[dv=dom1]    (oth) at (10.4,-0.4){\faIcon{server}\\Everything\\else};
  \draw[w] (h) -- (r1);  \draw[w] (r1) -- (r2);
  \draw[w] (r2) -- (fin);  \draw[w] (r2) -- (oth);

  % 2 -- the obvious place, and the outage
  \stepvis{2}{%
    \draw[accent!85, line width=1.4pt] (4.7,1.25) -- (5.3,0.15);
    \node[font=\tiny\bfseries, text=accent!75!black, align=center,
          text width=3.1cm, anchor=north] at (5.0,-0.35)
        {\faIcon{times-circle}~\textbf{standard} here\\blocks
         \emph{everything}};}

  % 3 -- move it
  \stepvis{3}{%
    \draw[dom2!80, line width=1.4pt] (8.5,2.0) -- (9.1,1.1);
    \node[font=\tiny\bfseries, text=dom2!45!black, align=center,
          text width=3.4cm, anchor=south] at (8.8,2.15)
        {\faIcon{check-circle}~\textbf{standard} here\\blocks only this
         server};}

  \steponly{1}{\node[stepcap, text=neutral!45!black] at (-1.5,-1.6)
      {Two servers, one path. Whatever we write has to separate them.};}
  \steponly{2}{\node[stepcap, text=accent!75!black] at (-1.5,-1.6)
      {Near the source looks right --- it stops doomed traffic early. But a
       standard ACL matches only \emph{who sent it}, and here both servers
       are behind the same next hop. Everything VLAN 30 sends is dropped.};}
  \steponly{3}{\node[stepcap, text=dom2!45!black] at (-1.5,-1.6)
      {Move it to the last hop and the two servers are finally on different
       interfaces, so matching the source alone is enough.};}
  \steponly{4}{\node[stepcap, text=dom3!45!black] at (-1.5,-1.6)
      {Hence the rule: a \textbf{standard} ACL goes near the destination,
       because that is the only place its one match is sufficient. An
       \textbf{extended} ACL names the destination too, so it is safe at the
       source --- and saves carrying doomed traffic across the network.};}
\end{tikzpicture}
